\documentclass[12pt,a4paper]{article}
\usepackage[utf8]{inputenc}
\usepackage[french]{babel}
\usepackage[T1]{fontenc}
\usepackage{geometry}
\geometry{margin=2.5cm, top=3cm}
\usepackage{xcolor}
\usepackage{titlesec}
\usepackage{listings}
\usepackage{tikz}
\usepackage[most]{tcolorbox}
\usepackage{inconsolata}
\usepackage{textcomp}

% --- Palette de Couleurs ---
\definecolor{MainBlue}{RGB}{20, 60, 110}
\definecolor{LightBlue}{RGB}{235, 245, 255}
\definecolor{TerminalBlack}{RGB}{40, 44, 52}
\definecolor{ShadowGray}{RGB}{210, 210, 210}
\definecolor{CodeKeyword}{RGB}{181, 137, 0}
\definecolor{CodeString}{RGB}{42, 161, 152}
\definecolor{CodeComment}{RGB}{133, 153, 0}

% --- Section Title Styling ---
\titleformat{\section}
{\color{MainBlue}\normalfont\Large\bfseries}
{}{0em}
{%
  \begin{tikzpicture}[baseline=(char.base)]
    \fill[MainBlue] (0,-0.1) rectangle (0.3,0.6);
    \node (char) at (0.5,0) {};
  \end{tikzpicture}\hspace{0pt}%
}

% ---- Java/JSP listing style ----
\lstdefinestyle{pystyle}{
    language=Java,
    basicstyle=\ttfamily\footnotesize,
    keywordstyle=\color{CodeKeyword}\bfseries,
    stringstyle=\color{CodeString},
    commentstyle=\color{CodeComment}\itshape,
    numberstyle=\ttfamily\tiny\color{gray!60},
    columns=fixed,
    keepspaces=true,
    upquote=true,
    showstringspaces=false,
    numbers=left,
    numbersep=10pt,
    tabsize=4,
    breaklines=true,
    breakatwhitespace=false,
    breakindent=2em,
    postbreak=\mbox{\textcolor{gray}{\ensuremath{\hookrightarrow}}\space},
}

% ---- Code Box ----
\newtcblisting{codeholder}[1]{
    listing only,
    listing style=pystyle,
    enhanced,
    breakable,
    colback=white,
    colframe=TerminalBlack,
    boxrule=0.7pt,
    arc=4pt,
    left=2pt,
    right=4pt,
    top=4pt,
    bottom=4pt,
    boxsep=4pt,
    title={%
        \textcolor{red!80}{$\bullet$}
        \textcolor{yellow!80}{$\bullet$}
        \textcolor{green!80}{$\bullet$}
        \hspace{0.4cm}
        {\ttfamily\footnotesize #1}
    },
    fonttitle=\bfseries\color{white},
    colbacktitle=TerminalBlack,
    attach boxed title to top left={yshift=-3mm, xshift=6mm},
    boxed title style={size=small, sharp corners, arc=2pt},
    shadow={1mm}{-1mm}{0mm}{ShadowGray},
}

\begin{document}

\begin{titlepage}
    \begin{tikzpicture}[remember picture, overlay]
        \fill[MainBlue] (current page.north west) rectangle ([xshift=1.5cm]current page.south west);
        \node[anchor=north west, fill=LightBlue, minimum width=15cm, minimum height=4cm, inner sep=1cm] at ([xshift=2cm, yshift=-4cm]current page.north west) {
            \begin{minipage}{14cm}
                {\Huge \bfseries \color{MainBlue} RAPPORT DE TP5} \\[0.4cm]
                {\large \color{black!70} Développement d'une Application Java EE}
            \end{minipage}
        };
    \end{tikzpicture}
    \vfill
    \begin{flushright}
    \begin{tcolorbox}[width=8.5cm, colframe=MainBlue, colback=white, boxrule=2pt, arc=10pt, inner sep=15pt, shadow={3mm}{-3mm}{0mm}{ShadowGray}]
        {\large \textbf{IDENTITÉ}} \vspace{5pt} \hrule \vspace{10pt}
        \begin{tabular}{ll}
            \textbf{Nom :} & Jalal Grini \\
            \textbf{Classe :} & TDIA2 \\
            \textbf{École :} & ENSAH \\
            \textbf{Enseignant :} & M.Cherradi \\
            \textbf{Sujet :} & Application Web CRUD \\ 
                            & MVC2 avec Hibernate
        \end{tabular}
    \end{tcolorbox}
    \end{flushright}
    \vspace{2cm}
\end{titlepage}

\newpage
\tableofcontents
\newpage

\section{Introduction et Architecture}
Ce projet a pour objectif de développer une application web de gestion de produits. L'application s'appuie sur une architecture robuste séparant clairement les responsabilités :
\begin{itemize}
    \item \textbf{Hibernate ORM} : Gère le mapping objet-relationnel avec MySQL.
    \item \textbf{La Logique Métier (DAO et Services)} : Gère les règles d'application et l'accès à la base de données.
    \item \textbf{Les Vues (JSP)} : Se contentent d'afficher les données préparées par les contrôleurs.
\end{itemize}

\section{Étape 1 : Configuration (Hibernate)}
\begin{codeholder}{src/main/resources/hibernate.cfg.xml}
<hibernate-configuration>
    <session-factory>
        <property name="hibernate.connection.url">jdbc:mysql://localhost:3306/tpnewmvc2</property>
        <property name="hibernate.connection.username">root</property>
        <property name="hibernate.connection.password">Jalal.1234</property>
        <property name="hibernate.hbm2ddl.auto">update</property>
        <mapping class="DAO.Produit"/>
        <mapping class="model.User"/>
    </session-factory>
</hibernate-configuration>
\end{codeholder}

\begin{codeholder}{src/main/java/util/HibernateUtil.java}
public class HibernateUtil {
    private static SessionFactory sessionFactory = null;
    public static SessionFactory getSessionFactory() {
        if (sessionFactory == null) {
            sessionFactory = new Configuration().configure("hibernate.cfg.xml").buildSessionFactory();
        }
        return sessionFactory;
    }
}
\end{codeholder}

\section{Étape 2 : Les Modèles de Données}
\begin{codeholder}{src/main/java/model/User.java}
@Entity @Table(name = "users")
public class User {
    @Id private String login;
    private String password;
    private String role;
    
    public User(String login, String password, String role) {
        this.login = login; this.password = password; this.role = role;
    }
}
\end{codeholder}

\begin{codeholder}{src/main/java/DAO/Produit.java}
@Entity @Table(name = "produits")
public class Produit {
    @Id @GeneratedValue(strategy = GenerationType.IDENTITY)
    private Long idProduit;
    private String nom;
    private String description;
    private Double prix;
}
\end{codeholder}

\section{Étape 3 : Accès aux Données (DAO)}
\begin{codeholder}{src/main/java/DAO/ProduitDAO.java}
public interface ProduitDAO {
    public void addProduit(Produit p);
    public void deleteProduit(Long id);
    public void updateProduit(Produit p);
    public List<Produit> getAllProduits();
    public Produit getProduitById(Long id);
}
\end{codeholder}

\begin{codeholder}{src/main/java/DAO/ProduitImpl.java}
public class ProduitImpl implements ProduitDAO {
    @Override
    public void addProduit(Produit p) {
        Session s = HibernateUtil.getSessionFactory().openSession();
        Transaction tx = s.beginTransaction();
        s.save(p); tx.commit(); s.close();
    }
    @Override
    public void deleteProduit(Long id) {
        Session s = HibernateUtil.getSessionFactory().openSession();
        Transaction tx = s.beginTransaction();
        Produit p = s.get(Produit.class, id);
        if (p != null) s.delete(p);
        tx.commit(); s.close();
    }
    @Override
    public List<Produit> getAllProduits() {
        Session s = HibernateUtil.getSessionFactory().openSession();
        List<Produit> list = s.createQuery("from Produit", Produit.class).list();
        s.close(); return list;
    }
    // Reste de l'implémentation (update, getById)...
}
\end{codeholder}

\begin{codeholder}{src/main/java/DAO/UserDAO.java}
public class UserDAO {
    public User findByLoginAndPassword(String login, String password) {
        Session session = HibernateUtil.getSessionFactory().openSession();
        User u = session.createQuery("from User where login = :l and password = :p", User.class)
                .setParameter("l", login).setParameter("p", password).uniqueResult();
        session.close(); return u;
    }
}
\end{codeholder}

\section{Étape 4 : Couche Métier (Services)}
\begin{codeholder}{src/main/java/Services/ProduitMetier.java}
public interface ProduitMetier {
    public void addProduit(Produit p);
    public void deleteProduit(Long id);
    public List<Produit> getAllProduits();
    public Produit getProduitById(Long id);
}
\end{codeholder}

\begin{codeholder}{src/main/java/Services/ProduitMetierImpl.java}
public class ProduitMetierImpl implements ProduitMetier {
    private static ProduitMetierImpl instance;
    private ProduitDAO dao = new ProduitImpl();
    public static ProduitMetierImpl getInstance() {
        if (instance == null) instance = new ProduitMetierImpl();
        return instance;
    }
    @Override
    public void addProduit(Produit p) { dao.addProduit(p); }
    @Override
    public List<Produit> getAllProduits() { return dao.getAllProduits(); }
}
\end{codeholder}

\section{Étape 5 : Sécurité (Filtre)}
\begin{codeholder}{src/main/java/filters/SessionFilter.java}
@WebFilter(urlPatterns = { "/listProduits", "/addProduit", "/deleteProduit", "/editProduit", "/updateProduit" })
public class SessionFilter extends HttpFilter {
    public void doFilter(ServletRequest req, ServletResponse resp, FilterChain chain) throws ... {
        HttpServletRequest request = (HttpServletRequest) req;
        HttpSession session = request.getSession(false);
        if (session != null && session.getAttribute("user") != null) {
            chain.doFilter(req, resp);
        } else {
            ((HttpServletResponse)resp).sendRedirect(request.getContextPath() + "/login");
        }
    }
}
\end{codeholder}

\section{Étape 6 : Les Contrôleurs (Servlets)}
\begin{codeholder}{src/main/java/web/LoginServlet.java}
@WebServlet("/login")
public class LoginServlet extends HttpServlet {
    private UserDAO dao = new UserDAO();
    protected void doPost(HttpServletRequest request, HttpServletResponse response) throws ... {
        User u = dao.findByLoginAndPassword(request.getParameter("login"), request.getParameter("password"));
        if (u != null) {
            request.getSession().setAttribute("user", u);
            response.sendRedirect("listProduits");
        } else {
            request.setAttribute("erreur", "Invalide");
            request.getRequestDispatcher("login.jsp").forward(request, response);
        }
    }
}
\end{codeholder}

\begin{codeholder}{src/main/java/web/ListProduitsServlet.java}
@WebServlet("/listProduits")
public class ListProduitsServlet extends HttpServlet {
    protected void doGet(HttpServletRequest request, HttpServletResponse response) throws ... {
        request.setAttribute("listeProduits", ProduitMetierImpl.getInstance().getAllProduits());
        User u = (User) request.getSession().getAttribute("user");
        String page = ("admin".equals(u.getRole())) ? "index.jsp" : "list.jsp";
        request.getRequestDispatcher(page).forward(request, response);
    }
}
\end{codeholder}

\begin{codeholder}{src/main/java/web/AddProduitServlet.java}
@WebServlet("/addProduit")
public class AddProduitServlet extends HttpServlet {
    protected void doPost(HttpServletRequest req, HttpServletResponse resp) throws ... {
        ProduitMetierImpl.getInstance().addProduit(new Produit(
            req.getParameter("nom"), req.getParameter("description"), 
            Double.parseDouble(req.getParameter("prix"))
        ));
        resp.sendRedirect("listProduits");
    }
}
\end{codeholder}

\begin{codeholder}{src/main/java/web/EditProduitServlet.java}
@WebServlet("/editProduit")
public class EditProduitServlet extends HttpServlet {
    protected void doGet(HttpServletRequest req, HttpServletResponse resp) throws ... {
        Long id = Long.parseLong(req.getParameter("id"));
        req.setAttribute("produitEdit", ProduitMetierImpl.getInstance().getProduitById(id));
        req.setAttribute("listeProduits", ProduitMetierImpl.getInstance().getAllProduits());
        req.getRequestDispatcher("index.jsp").forward(req, resp);
    }
}
\end{codeholder}

\begin{codeholder}{src/main/java/web/UpdateProduitServlet.java}
@WebServlet("/updateProduit")
public class UpdateProduitServlet extends HttpServlet {
    protected void doPost(HttpServletRequest req, HttpServletResponse resp) throws ... {
        Produit p = new Produit(req.getParameter("nom"), req.getParameter("description"), Double.parseDouble(req.getParameter("prix")));
        p.setIdProduit(Long.parseLong(req.getParameter("idProduit")));
        ProduitMetierImpl.getInstance().updateProduit(p);
        resp.sendRedirect("listProduits");
    }
}
\end{codeholder}

\begin{codeholder}{src/main/java/web/DeleteProduitServlet.java}
@WebServlet("/deleteProduit")
public class DeleteProduitServlet extends HttpServlet {
    protected void doGet(HttpServletRequest req, HttpServletResponse resp) throws ... {
        ProduitMetierImpl.getInstance().deleteProduit(Long.parseLong(req.getParameter("id")));
        resp.sendRedirect("listProduits");
    }
}
\end{codeholder}

\begin{codeholder}{src/main/java/web/LogoutServlet.java}
@WebServlet("/logout")
public class LogoutServlet extends HttpServlet {
    protected void doGet(HttpServletRequest req, HttpServletResponse resp) throws ... {
        HttpSession session = req.getSession(false);
        if (session != null) session.invalidate();
        resp.sendRedirect("login");
    }
}
\end{codeholder}

\section{Étape 7 : Présentation (Vues JSP)}
\begin{codeholder}{src/main/webapp/login.jsp}
<form action="login" method="post">
    Login: <input type="text" name="login" required /><br/>
    Password: <input type="password" name="password" required /><br/>
    <input type="submit" value="Se connecter" />
</form>
\end{codeholder}

\begin{codeholder}{src/main/webapp/index.jsp}
<form action="${produitEdit != null ? 'updateProduit' : 'addProduit'}" method="post">
    <input type="hidden" name="idProduit" value="${produitEdit.idProduit}" />
    Nom: <input type="text" name="nom" value="${produitEdit.nom}" required />
    Prix: <input type="text" name="prix" value="${produitEdit.prix}" required />
    <input type="submit" value="${produitEdit != null ? 'Modifier' : 'Ajouter'}" />
</form>
<table border="1">
    <c:forEach var="p" items="${listeProduits}">
        <tr>
            <td>${p.idProduit}</td><td>${p.nom}</td><td>${p.prix}</td>
            <td><a href="editProduit?id=${p.idProduit}">Edit</a> | <a href="deleteProduit?id=${p.idProduit}">Del</a></td>
        </tr>
    </c:forEach>
</table>
\end{codeholder}

\begin{codeholder}{src/main/webapp/list.jsp}
<table border="1">
    <tr><th>ID</th><th>Nom</th><th>Prix</th></tr>
    <c:forEach var="p" items="${listeProduits}">
        <tr><td>${p.idProduit}</td><td>${p.nom}</td><td>${p.prix}</td></tr>
    </c:forEach>
</table>
\end{codeholder}

\begin{codeholder}{src/main/webapp/error.jsp}
<h2>Erreur</h2>
<p style="color:red;"><%= exception.getMessage() %></p>
<a href="listProduits">Retour</a>
\end{codeholder}

\end{document}
